\section{Introduction}
Railway networks these days face continuous pressure to optimize operations and reduce delays, while also operating energy efficiently. Within the domain of railway scheduling, finding conflict-free routes for multiple agents on a shared grid is a complex combinatorial challenge. To study these dynamics, researchers often use the Flatland environment, which serves as a benchmarking tool for multi-agent coordination. While much of the existing research uses Reinforcement Learning to navigate Flatland's grid worlds, our project adopts a declarative approach by utilizing Answer Set Programming (ASP) to formulate the scheduling task as a series of logical constraints, allowing the solver to rigorously explore the combinatorial space for optimal routing.

Traditional pathfinding in Flatland generally focuses on one primary objective: navigating agents from a starting point to a destination in the shortest amount of time. However, real-world railway logistics demand the optimization of secondary, physically grounded metrics. In this paper, we extend Flatland's routing logic by introducing energy efficiency as a core constraint. We model the physical realities of train operations by calculating energy consumption dynamically based on individualized vehicle mass, motor efficiency (electric, diesel, hybrid, and hydrogen), and kinematic states such as waiting, moving straight ahead and navigating turns. By treating railway scheduling as a multi-objective optimization problem, we aim to evaluate how the introduction of energy penalties influences the routing choices made by the ASP solver. 

The remainder of this report is structured as follows. In Section 2, we introduce the baseline ASP encoding required for basic vehicle routing and collision avoidance within Flatland. We then detail the mechanics of our energy encoding, explaining how physical traits and movement penalties are mathematically integrated into the logic. In Section 3, we present the results of our experiments, analyzing how the solver navigates the trade-offs between arrival time and fuel efficiency. Finally, in Section 4, we conclude on our findings regarding the computational complexity of the encoding and suggest future directions for research in multi-objective railway scheduling.